%=============================================================================
% Mechanism Is Not Meaning:
% A Verification Methodology for Cognitive-Architecture Prototypes
%
% STRUCTURE-ONLY SKELETON. No body prose. Each \section carries a single
% planning comment describing its future content. The body is written later,
% bound by the locked claim inventory that governs every claim.
%
% PERMANENT CONSTRAINTS (whole paper):
%  - Self-contained: every mechanism, algorithm, and rationale is explained in
%    the paper text itself. Defer nothing to external code or files.
%  - Zero source references: no commit hashes, no file paths, no external URL,
%    and no external-location pointer. Claims rest on the in-text mechanisms and
%    on the evidence in the appendices (the underlying project stays private).
%  - Appendices carry key evidence: clean tables as the main vehicle, plus a
%    few verbatim raw-output excerpts at the most critical points.
%  - No em-dashes (never three consecutive hyphens); en-dash (two hyphens) only
%    for numeric ranges. Demonstrate rigor through calibrated, evidence-bound
%    claims; certainty matched to each claim's assurance tier; full
%    publication-quality academic register.
%  - System name in the paper is "Psyche"; the codebase stays "Nous".
%=============================================================================
\documentclass[11pt]{article}

% Suppress one benign cosmetic warning from lmodern ("Command \showhyphens
% has changed") seen on current TeX Live. Guarded so the document still
% compiles if the silence package is unavailable; must precede lmodern.
\IfFileExists{silence.sty}{%
  \usepackage{silence}%
  \WarningFilter{lmodern}{Command}%
}{}

% Input encoding and fonts (Turkish glyphs in the author name)
\usepackage[utf8]{inputenc}
\usepackage[T1]{fontenc}
\usepackage{lmodern}

% Layout and mathematics
\usepackage[margin=1in]{geometry}
\usepackage{amsmath}

% Graphics, tables, diagrams
\usepackage{graphicx}
\usepackage{booktabs}
\usepackage{tikz}

% Bibliography (author-year via natbib + bibtex)
\usepackage{natbib}
\bibliographystyle{plainnat}

% Hyperlinks (loaded after most packages; xcolor arrives via tikz)
\usepackage{hyperref}
\hypersetup{
  colorlinks = true,
  linkcolor  = black,
  citecolor  = blue!45!black,
  urlcolor   = blue!45!black,
  pdftitle   = {Mechanism Is Not Meaning: A Verification Methodology for Cognitive-Architecture Prototypes},
  pdfauthor  = {Cagan Caliskan}
}

% ORCID: use orcidlink if installed, otherwise degrade gracefully
\IfFileExists{orcidlink.sty}{%
  \usepackage{orcidlink}%
  \newcommand{\authororcidicon}{\,\orcidlink{0009-0000-2993-1421}}%
}{%
  \newcommand{\authororcidicon}{}%
}

% Microtypographic refinement when available
\IfFileExists{microtype.sty}{\usepackage{microtype}}{}

%=============================================================================
\title{Mechanism Is Not Meaning:\\[2pt]
       A Verification Methodology for Cognitive-Architecture Prototypes}

\author{%
  Çağan Çalışkan\authororcidicon \\[4pt]
  \normalsize ORCID:~\href{https://orcid.org/0009-0000-2993-1421}{0009-0000-2993-1421}%
}

\date{June 2026}

%=============================================================================
\begin{document}
\maketitle

\begin{abstract}
% ABSTRACT - to be written last
\end{abstract}

%=============================================================================
\section{Introduction}
% Problem: separating mechanism from meaning; false-positive cognitive claims;
% Psyche as the case study; the methodology is THE central contribution; one
% sentence on the long-term non-LLM motivation, with scope explicitly limited
% and not evaluated here.
% NOTE: the first occurrence of "Psyche" in this section carries the branding
% bridge footnote, used exactly ONCE:
%   \footnote{Psyche (developed in the repository under the working name Nous).}

%=============================================================================
\section{The Substrate and the Closed Cognitive Cycle}

Psyche is built in two layers that carry two different grades of evidence, and the
distinction between them is maintained throughout this paper. The lower layer, the
\emph{substrate}, is a graph store together with a spreading-activation engine; its
properties are proven, in the precise sense established below: identical output
across separate processes under randomized hashing, mutual consistency of the graph
indices after every mutation, and reads that never alter stored state. The upper
layer assembles the substrate into a single loop, perceiving, activating,
clustering, selecting, pruning, and writing, whose stages are verified as mechanisms
but make no claim to reproduce a cognitive phenomenon. We call the loop a
\emph{cognitive cycle} because its shape is taken from the global-workspace tradition
\citep{baars1988cognitive, franklin2014lida}, as a structural template rather than a
theory the system is claimed to instantiate; the term names a control structure, and
its use here asserts nothing about cognition. What the loop is shown to do is to
\emph{close}, as a mechanism, and that is a narrower and more exact statement than it
may first appear.

\subsection{The substrate store}

The store is a directed graph whose vertices are concept \emph{nodes} and whose edges
are weighted \emph{roads}. Nodes hold no state of their own; everything that varies
lives on the roads. A road carries a scalar \emph{strength} together with a record of
how it has been used, advances through a small set of discrete quality tiers as it is
traversed, and decays with elapsed time when it is not. A road may also be
\emph{closed}: a closed road remains in the graph and stays traversable by a
sufficiently strong cue, so it is leaky rather than severed, but it no longer conducts
ordinary activation.

The load-bearing decision in the store is that reading and writing are separated at
the level of decay. An earlier arrangement applied decay at the moment of access and
wrote the decayed value back, with two consequences a casual reading would miss:
every read mutated the store, and, because that same step re-evaluated closure, a
sequence of reads spread over time could let a closed road drift back open although
nothing had ever crossed it. Psyche instead splits decay into a pure view and a
persisting write. A read returns a decayed \emph{copy}, computed from the elapsed time
but stored nowhere, so reading changes nothing; only the active operations,
traversal and reinforcement and the offline maintenance passes, persist a decayed
value. Closure is made \emph{sticky} by the same principle: it is re-evaluated only on
an active qualifying crossing, never by elapsed time and never by a read, so no amount
of passive inspection can reopen what only an event should reopen. These are the
properties the substrate is proven to hold: reads are pure, closure does not drift,
the graph indices stay mutually consistent through every insertion and deletion
(including on adversarially constructed graphs), and a deletion removes a node
completely, leaving no dangling entry in any index. The corresponding assertion counts
appear in Appendix~\ref{app:verification}.

\subsection{The spreading-activation engine}

Activation is the substrate's one dynamic operation. Given a set of seed nodes, the
engine propagates a diminishing quantity outward along the roads and reports how much
activation each reached node accumulates. The mechanism is the spreading activation of
ACT-R \citep{anderson1998atomic}; Psyche borrows it without that architecture's
empirical commitments. Two design choices make the
computation well defined.

The first is bounded single-expansion. When activation arriving along several paths
converges on a node, the quantities must add; the hazard is that the node, once
re-expanded as later paths raise its value, propagates the surplus, and everything
downstream of the convergence point is counted more than once. The error is the
residue of a shortest-path intuition, in which the first time a node is settled is
final, carried into a setting where contributions accumulate rather than compete.
Psyche expands each node exactly once, at the first hop at which it is reached,
carrying the full sum of all paths of length up to that hop; activation that converges
by equal-length paths is therefore summed once and exactly, independent of the
relative strengths of the converging roads.

The second choice secures determinism. If expansion follows an activation-ordered
frontier, which late contribution is discarded depends on incidental orderings of edge
strength. Separately, if same-hop siblings expand in the order a set of
machine-generated identifiers happens to iterate, that order is randomized from one
process to the next by the hashing of those identifiers, and the result moves between
runs for no reason present in the graph. Psyche instead expands hop by hop and freezes each layer before
it acts: all contributions that reach distance $k$ are accumulated into a snapshot,
the snapshot is frozen, and only then do the hop-$k$ nodes expand, reading their values
from the frozen snapshot rather than from one another. Because no node within a layer
can observe another's expansion, the outcome cannot depend on the order in which nodes
or their outgoing roads are visited. The engine is in consequence deterministic across
separate processes under randomized process hash seeds (5--8 per suite, in fresh
interpreters), and, like the store, it computes as a pure read: a spread mutates
nothing.

Together the store and the engine form the substrate, the proven layer of the system.
Their two suites contribute 45 and 35 assertions respectively, 80 in all, within a
verification suite that currently stands at 20 suites and 1,372 assertions with no
failures (Appendix~\ref{app:verification}). Everything above this layer is built on it
and claims less.

\subsection{The closed cognitive cycle}

Above the substrate, Psyche arranges the operations of the store and the engine into a
single loop, shown in Figure~\ref{fig:cycle}. A cycle begins by \emph{perceiving}: the
input text is matched against the substrate, and the nodes it names become the cycle's
sources. This step is what makes the rest of the loop a genuine test rather than a staged one,
because the sources are produced by the system from its own contents rather than placed
by hand; the loop is driven by input it generates, not by input chosen to make it
succeed.

\begin{figure}[t]
  \centering
  % TikZ F-1 to be drawn: the closed cognitive cycle as information flow.
  % Stages, in order: perceive, activate, cluster, select, prune, write (broadcast),
  % then back to the next cycle. Use neutral verbs only (select, prune, write);
  % do not use decide, attend, or experience. Mark the feedback edge from the write
  % stage to the next cycle as "bounded, default-off feedback".
  \caption{The closed cycle of Psyche as information flow at the level of mechanism. A
    cycle proceeds from perception to activation, clusters the perceived sources,
    selects a single winning coalition, prunes the region that coalition reaches across
    a bridge, and performs one bounded write, the broadcast, that biases the next
    cycle. The figure depicts graph operations and a single bounded write; the verbs
    are select, prune, and write, not decide, attend, or experience. The feedback edge
    is bounded, default-off feedback. No cognitive process is depicted, and no
    cognitive claim is made.}
  \label{fig:cycle}
\end{figure}

Perception rarely yields a clean set of well-separated sources; it names several that
lie close together in the graph, and an activation field built directly on such a set
produces basins that overlap and collapse rather than separate concerns. The first
stage after perception therefore groups the sources before any field is constructed.
Psyche builds a small proximity graph over the perceived sources, using only the
strength of the roads between them, merges them by single linkage, and seeds the
downstream field from one representative per group. The rationale is that the grouping
must come from the structure of the substrate, not from the activation it is about to
compute, so that the later stages are not handed a partition that already encodes the
answer; using road strength alone keeps the pre-stage topological. The merge is a
deterministic, pure, integrity-preserving operation, and it is verified as such: it
asserts nothing about meaning, only about which sources the graph places together. The
threshold that governs how close is close enough is a fixed design constant rather than
a calibrated one, and we mark it as such instead of presenting it as tuned.

From the grouped sources the engine computes its field, and the co-activated region is
partitioned into candidate coalitions, of which exactly one is selected as the cycle's
result. Selection is a competition normalized for coalition size, so that a coalition
does not prevail merely by being large, and it returns a single winner
deterministically. Selection then prunes from the winner any region it has reached
across a bridge: a perceived concern can extend, through a single connecting road, into
a neighbouring region the input never named, and that unnamed region should not enter
the result on the strength of one bridge. The pruning rule is parameter-free. A reached
element is kept only if it is attached to the named region at least as strongly as to
the unnamed one; where attachment is stronger toward the unnamed side, the rule finds a
cohesion valley, a seam of relatively weak connection, and cuts there. The rule has one
corner case, which we state rather than conceal: an unnamed region bridged into the
named one at least as strongly as it coheres internally presents no valley to cut at,
and is retained. The divergence is one-directional and on the safe side, in that the
rule may keep an unnamed region a stricter criterion would drop but never prunes a
member the input actually named. Selection, pruning included, is a pure read; it is
verified as a mechanism and asserts nothing about cognition.

Up to this point the cycle has only read the substrate. Its one write is the final
stage, the broadcast: the winning coalition is fed back onto the substrate by
reinforcing its internal roads, biasing the next cycle toward the same region. This
single write is what turns a pipeline into a loop, and it is also the first opportunity
the substrate has to corrupt itself over repeated cycles, so its boundedness is
established in both directions. The reinforcement reuses the store's ordinary
strengthening rule, which adds to a road's strength a fixed fraction of the distance
that remains to its ceiling,
\[
  s \;\leftarrow\; s + \gamma\,(1 - s), \qquad 0 < \gamma \le 1,
\]
so that each application closes part of the gap to $s = 1$ without ever reaching or
crossing it. Boundedness is thus a structural consequence of the update rather than a
limit imposed afterwards: the broadcast contains no cap, no ceiling test, and no
tunable magnitude of its own, the gain $\gamma$ being the store's existing default. In
the other direction, the feedback does not lock the loop onto its first winner; a
sufficiently strong later input still moves the result, which was confirmed by driving
one region to dominance over sustained reinforcement and then verifying that a
competing input takes the workspace, with the first region checked to be dominant
immediately before it yields. The broadcast is off by default, and with it off the loop
is a pure read whose output is byte-for-byte identical to the pre-broadcast pipeline;
enabling it adds exactly one bounded write per cycle and nothing else.

With the broadcast in place the loop is closed: perception, activation, clustering,
selection, pruning, and the single feedback write run end to end on input the system
generates, the read portion is deterministic across processes and pure, and the write
portion is the one bounded reinforcement just described. This is the precise content of
the claim that the cycle closes. It is a statement about a mechanism, a pipeline that
runs to completion on its own output, repeatably and with no side effect beyond one
bounded write. It is not a statement that the loop perceives, attends, decides, or
means anything. Each of its stages, the clustering, the competition and pruning, and
the bounded broadcast, has been independently and adversarially re-verified for the
same properties the substrate is held to, namely determinism across randomized process
hash seeds, preserved integrity, and read-purity, and for none of them is any cognitive
claim entered; their assertion counts are recorded in Appendix~\ref{app:verification}.
The loop is built and its mechanism is verified. Whether anything it computes
corresponds to a cognitive phenomenon is a separate question, taken up, and left open,
later in this paper.

%=============================================================================
\section{The Verification Methodology}
% THE central contribution: four-agent separation; loop-closure; adversarial
% and property-based testing; pre-registration; anti-circularity gates.
% Figures F-2 (four-agent verification topology) and F-6 (six-bug table).
% Suite, assertion, and seed counts are given in-text and in
% Appendix~\ref{app:verification}; no external artifact is cited.

%=============================================================================
\section{The Methodology in Action: Three Meaning-Bar Attempts}
% Empirical evidence: 1e priming (circular), the fan effect (mechanism absent),
% 2c interoception (signal flat); each a NOT-YET with no graded number
% manufactured. Figures F-3a, F-3b, F-3c.
% The concrete numbers for all three are in Appendix~\ref{app:attempts-data};
% no external artifact is cited.

%=============================================================================
\section{The Meta-Finding}
% Three structural walls, one root: the external yardstick.
% Figures F-5 (meta-finding schematic) and F-4 (three-tier assurance map).

%=============================================================================
\section{Limitations}
% L-1 through L-8, stated plainly and up front.

%=============================================================================
\section{Related Work}
% Design lineage, NOT comparison; the dependency-free choice (why not QuickCheck
% or Hypothesis) belongs here.

%=============================================================================
\section{Conclusion}
% What was learned; what is claimed (methodology and mechanism) and what is not
% (no mind, consciousness, AGI, or validated cognition).

%=============================================================================
\appendix

%=============================================================================
\section{Verification Summary}\label{app:verification}
% Appendix A: a clean table of the canonical verification run (20 suites,
% 1,372 assertions, 0 failures), a per-suite breakdown by component, and the
% per-suite randomized process-hash-seed counts (5 to 8). This is the evidence
% that the mechanisms are verified; the numbers come from the canonical run.

%=============================================================================
\section{The Three Meaning-Bar Attempts: Data}\label{app:attempts-data}
% Appendix B: the concrete data behind the three NOT-YETs.
%  - 1e priming: N_S0 = 7, N_S+ = 296, the in-degree decile cuts
%    (1,1,1,1,1,2,2,5,15), and the unreachable-fallback arithmetic
%    (r_W ~ 0.97, so the needed r_MA exceeds 1.0).
%  - Fan effect: the flat sweep (target activation 0.2433743975, spread 0.0
%    across fan in {1,2,4,8,16}); include one short verbatim raw-output excerpt
%    here, the byte-identical-across-fan output being the clearest evidence.
%  - 2c interoception: the friction band, the populability pilot
%    (34 of 40 winner-flips, 85 per cent), and the scale-invariance argument.
% Clean tables are the main vehicle, plus the single fan-effect raw excerpt.

%=============================================================================
% Scaffolding only: \nocite{*} forces every bibliography entry to appear so the
% reference list can be checked now. Remove once real \cite commands populate the
% prose (final state: every \cite resolves, and every entry is cited).
\nocite{*}
\bibliography{references}

\end{document}
